\chapter{大型代码组织和演进}

\section{问题入口：一个 Manager 类吞掉项目}

很多项目会出现一个 \code{AppManager}：它读配置、连数据库、处理请求、写日志、管理缓存、启动线程。短期看方便，长期看每个修改都要碰它，测试也越来越难。

架构不是画复杂图，而是控制依赖方向。先把系统拆成几个角色：

\begin{itemize}
  \item \textbf{配置}：读取和验证启动参数。
  \item \textbf{领域逻辑}：处理核心业务，不知道文件和网络。
  \item \textbf{适配器}：把文件、数据库、HTTP 等外部世界接进来。
  \item \textbf{入口}：组装对象，建立异常和日志边界。
\end{itemize}

\section{依赖接口而不是具体环境}

假设业务逻辑需要读取用户信息。不要让它直接依赖文件路径：

\begin{lstlisting}[language=C++,caption={用接口隔离外部数据源}]
#include <optional>
#include <string>

struct UserRecord {
    int id = 0;
    std::string name;
};

class UserRepository {
public:
    virtual ~UserRepository() = default;
    [[nodiscard]] virtual std::optional<UserRecord> find_by_id(int id) const = 0;
};
\end{lstlisting}

业务服务依赖这个接口：

\begin{lstlisting}[language=C++,caption={业务逻辑依赖抽象}]
class GreetingService {
public:
    explicit GreetingService(const UserRepository& repository)
        : repository_{repository}
    {
    }

    [[nodiscard]] std::string greeting_for(int id) const
    {
        const std::optional<UserRecord> user = this->repository_.find_by_id(id);
        if (!user.has_value()) {
            return "unknown user";
        }
        return "hello, " + user->name;
    }

private:
    const UserRepository& repository_;
};
\end{lstlisting}

测试时可以传入内存实现；生产环境可以传入文件或数据库实现。

\section{继承要克制}

上面的接口用了虚函数，是为了隔离外部数据源。不要把所有变化点都做成继承。很多时候，普通值对象、函数对象、模板参数或构造函数注入更简单。

\begin{lstlisting}[language=C++,caption={小策略可以用函数对象}]
using PriceRule = std::function<int(int base_price)>;

int final_price(int base_price, const PriceRule& rule)
{
    return rule(base_price);
}
\end{lstlisting}

如果策略在热路径，\code{std::function} 的类型擦除成本可能需要测量；也可以用模板参数。架构决策要和性能约束一起看。

\section{模块边界和头文件}

头文件是依赖传播的入口。一个头文件包含太多其他头文件，会让编译变慢，也会放大耦合。优先在头文件里暴露必要类型，在源文件里包含实现细节。

\begin{lstlisting}[language=C++,caption={头文件保持小而稳定}]
#pragma once

#include <optional>
#include <string>

struct UserRecord;

class UserRepository {
public:
    virtual ~UserRepository() = default;
    [[nodiscard]] virtual std::optional<UserRecord> find_by_id(int id) const = 0;
};
\end{lstlisting}

不是所有类型都能前置声明，例如按值成员需要完整类型。知道这个边界，才能合理拆分。

\section{演进优先级}

当项目开始变大，优先处理这些问题：

\begin{enumerate}
  \item 重复业务规则：抽成清晰函数或类型。
  \item 隐式资源所有权：改成 RAII 和明确所有者。
  \item 巨大函数：按阶段拆分。
  \item 巨大类：按角色拆分。
  \item 难测试外部依赖：通过接口或适配器隔离。
\end{enumerate}

\begin{keyidea}
高级 \cpp 不是把模板、继承和并发全用上，而是在复杂系统里让依赖方向、资源生命周期和错误边界仍然清楚。
\end{keyidea}

\section{从难测试代码推导依赖方向}

如果 \code{GreetingService} 自己打开文件、解析数据、拼接问候语，测试一次问候语就要准备文件路径和文件内容。难测试不是测试人员的问题，而是依赖方向暴露了设计问题：业务逻辑依赖了外部环境。

把依赖倒过来后，业务逻辑只依赖 \code{UserRepository} 抽象。测试可以写一个内存仓库：

\begin{lstlisting}[language=C++,caption={测试用内存仓库}]
class InMemoryUserRepository final : public UserRepository {
public:
    void add(UserRecord user)
    {
        this->users_[user.id] = std::move(user);
    }

    [[nodiscard]] std::optional<UserRecord> find_by_id(int id) const override
    {
        const auto found = this->users_.find(id);
        if (found == this->users_.end()) {
            return std::nullopt;
        }
        return found->second;
    }

private:
    std::unordered_map<int, UserRecord> users_;
};
\end{lstlisting}

这段测试实现不需要文件和数据库。它也逼着接口保持小：如果 \code{UserRepository} 暴露十几个无关方法，测试替身就会很重。

\section{入口负责组装，不负责业务}

程序入口可以理解成装配层：

\begin{lstlisting}[language=C++,caption={入口组装真实依赖}]
int run_application(const Options& options)
{
    FileUserRepository repository{options.user_file};
    GreetingService service{repository};
    std::cout << service.greeting_for(options.user_id) << '\n';
    return 0;
}
\end{lstlisting}

这里可以出现文件路径、命令行选项、日志和异常边界，因为入口本来就在系统边缘。业务服务内部则不应该知道这些环境细节。这个边界能让核心代码在没有真实环境的情况下测试。

\section{什么时候不用虚函数}

虚函数适合运行期替换实现，例如生产环境和测试环境用不同仓库。如果策略在编译期就确定，模板或函数对象更轻：

\begin{lstlisting}[language=C++,caption={编译期策略组合}]
template <typename DiscountRule>
int final_price(int base_price, DiscountRule rule)
{
    return rule(base_price);
}

const int price = final_price(1000, [](int base) {
    return base - base / 10;
});
\end{lstlisting}

选择虚函数、模板还是 \code{std::function}，要看变化发生在运行期还是编译期、是否需要类型擦除、是否在热路径、报错和编译时间是否可接受。高级设计不是某个模式永远正确，而是能说出取舍。

\section{演进时不要大爆炸重写}

大型代码最怕“一次性重写”。更稳的演进方式是先找最痛的边界，建立薄接口，然后逐步迁移调用者。例如一个巨大 \code{AppManager} 可以先抽出配置读取：

\begin{lstlisting}[language=C++,caption={先抽出一个窄能力}]
class ConfigLoader {
public:
    [[nodiscard]] Config load(const std::string& path) const;
};
\end{lstlisting}

下一步再抽日志、缓存、业务服务。每抽一步都要保留测试，让行为不漂移。架构改进不是为了目录好看，而是为了让下一次修改更小、更可验证。

\begin{exercisebox}
选一个你写过的单文件程序，画出三类代码：输入输出、核心计算、格式化展示。先只抽出核心计算函数并写测试，不要急着引入类。等函数边界稳定后，再判断是否需要对象封装。
\end{exercisebox}
